\documentclass[10pt,a4paper]{amsart}
\usepackage[margin=19mm]{geometry}
\usepackage{amsmath,amssymb,mathtools}
\usepackage[scr=rsfs,cal=euler]{mathalfa}
\usepackage{xfrac}
\usepackage[unicode,hidelinks,bookmarks=false]{hyperref}
\newcommand{\ie}{i.e.\ }
\newcommand{\cf}{cf.\ }
\newcommand{\resp}{respectively\ }
\newcommand{\Subsection}{\subsection}
\providecommand{\nobreakdash}{\nobreak\hbox{-}}
\setlength{\emergencystretch}{2em}
\multlinegap0pt
\usepackage{booktabs}
\usepackage{nicefrac}
\usepackage{xcolor}
\usepackage{multirow}
\usepackage{mathtools}
\usepackage{algorithm}\usepackage{algpseudocode}
\newtheorem{ass}{Assumption}
\newtheorem{lem}{Lemma}
\newtheorem{prop}{Proposition}
\title{Tuning the burn-in phase in training recurrent neural networks improves their performance}
\author{Julian D. Schiller, Malte Heinrich, Victor G. Lopez, Matthias A. M\"uller}
\date{Source: arXiv:2602.10911v1 (CC BY 4.0)}
\begin{document}
\maketitle

\noindent\emph{Source and licence.} Tuning the burn-in phase in training recurrent neural networks improves their performance. Version arXiv:2602.10911v1 (CC BY 4.0), distributed by arXiv under the Creative Commons Attribution 4.0 International licence, http://creativecommons.org/licenses/by/4.0/, as recorded in the arXiv licence metadata for this exact version and shown on the abstract page https://arxiv.org/abs/2602.10911v1. Full licence notice: \texttt{licenses/arxiv-2602.10911-CC-BY-4.0.txt}.\par\smallskip

\noindent\textit{Editorial scope.} Counted unit: the article's prop prop:TP (source0.tex lines 654--661). The article's complete original proof of it is delivered with it (source0.tex lines 663--669, one \texttt{proof} environment of 7 lines). No mathematics is written, repaired or re-derived: the statement and the proof are the article's own lines, in the article's own notation, and no retained line is substituted or re-flowed.  Retained uncounted as the source-local context the proof needs: lines 111--113; lines 116--118; lines 281--287; lines 554--560; lines 590--592; lines 624--629.  Nothing was omitted from any retained span, so the package carries no omission record.  No retained line contains a citation, so the wrapper carries no bibliography and no bibliography entry.  No retained line contains a figure environment, an image-inclusion command or drawing code, so no image is delivered and the package declares no asset dependency.  Editorial formatting only, in addition: an AMS \texttt{amsart} preamble that restates the article's own declarations and macros used by the retained lines, namely the environments \texttt{ass}, \texttt{lem}, \texttt{prop} on separate counters, together with the standard packages those lines need.  The retained environments print in this document as Ass 1, Lemma 1, Proposition 1 in order of appearance, whereas the article prints the same items with the numbering of its own full document.  The complete native source of the article as distributed in the arXiv e-print for this exact version is bundled unchanged as \texttt{sources/194442/source0.tex}.
\begin{equation}\label{eq:RNN_1}
	h_{t} = f(h_{t-1},x_t;\theta_\mathrm{h}), \quad t = 1,...,T,
\end{equation}

\begin{equation}\label{eq:RNN_2}
	y_t = g(h_t,x_t;\theta_\mathrm{y}), \quad t=1,...,T,
\end{equation}

\begin{ass}\label{ass:output_stability}
	There exists a (non-empty) set $\Theta\subset\mathbb{R}^{n}$ such that, for all $\theta\in\Theta$, the network~\eqref{eq:RNN_1}--\eqref{eq:RNN_2} is exponentially incrementally output stable along input sequences of the training data, that is, there exist constants $C>0$ and $\lambda\in(0,1)$ such that
	\begin{equation}
		\|y_t(h_0^{(1)},{\theta},X)-y_t(h_0^{(2)},{\theta},X)\| \leq 	C\lambda^t\|h_0^{(1)}-h_0^{(2)}\|,\ t=1,2,...,T' \label{eq:output_stability}
	\end{equation}
	for any two initial hidden states $h^{(1)}_{0},h^{(2)}_{0}\in\mathbb{R}^n$ and any sequence $X$ of length $T'\in[0,T]$ extracted from the training input sequence $X^\mathrm{d}$.
\end{ass}

\begin{lem}[Stabilizability]\label{lem:controllability}
	Let Assumption~\ref{ass:output_stability} hold. Then, there exists a constant $\bar{C}>0$ such that
	\begin{equation}\label{eq:controllability}
		\sum_{i=1}^{S}\sum_{j=m+1}^{N}l_{j|i}({y}_{j|i}^*)-l_{j|i}({y}^\infty_{j|i}) \leq \bar{C}S\lambda^m,
	\end{equation}
	where $\lambda\in(0,1)$ is from Assumption~\ref{ass:output_stability}.
\end{lem}

	\begin{equation}\label{eq:lem_opt}
		\sum_{i=1}^{S}\sum_{j=m+1}^{N} 2({y}_{j|i}^\infty-{y}_{j|i}^\mathrm{d})^\top({y}_{j|i}^*-{y}_{j|i}^\infty) \geq-\frac{1}{\epsilon}\sum_{i=1}^{S}\sum_{j=m+1}^{N}\|{y}_{j|i}^*-{y}_{j|i}^\infty\|^2.
	\end{equation}

\begin{lem}[Coercivity]\label{lem:coercivity}
	Let the property in \eqref{eq:lem_opt} hold for some $\epsilon>1$. Then, the solutions $(H^*,Y^*,\theta^*)$ and $(H^\infty,Y^\infty,\theta^\infty)$ satisfy
	\begin{equation}\label{eq:lem_coercive}
		\sum_{i=1}^{S}\sum_{j=m+1}^{N}l_{j|i}({y}_{j|i}^*)-l_{j|i}({y}^\infty_{j|i}) \geq \frac{\epsilon-1}{\epsilon} \sum_{i=1}^{S}\sum_{j=m+1}^{N}\|{y}^*_{j|i}-{y}_{j|i}^\infty\|^2.
	\end{equation}
\end{lem}

\begin{prop}[Output turnpike property]\label{prop:TP}
	Let Assumption~\ref{ass:output_stability} be satisfied.
	Consider the solutions $(H^*,Y^*,\theta^*)$ and $(H^\infty,Y^\infty,\theta^\infty)$ and suppose that the property in \eqref{eq:lem_opt} holds for some $\epsilon>1$. Then, there exists a constant $K>0$ such that
	\begin{equation}\label{eq:prop_TP_1}
	 	\sum_{j=m+1}^{N} \ \frac{1}{S}\sum_{i=1}^{S}\|{y}_{j|i}^*-{y}_{j|i}^\infty\|^2 \leq K\lambda^m,
	\end{equation}
	where $\lambda\in(0,1)$ is from Assumption~\ref{ass:output_stability}.
\end{prop}

\begin{proof}
	The statement follows by considering the difference of the loss functions
	\begin{equation*}
		\sum_{i=1}^{S}\sum_{j=m+1}^{N}l_{j|i}({y}_{j|i}^*)-l_{j|i}({y}^\infty_{j|i})
	\end{equation*}
	 and applying the lower bound from Lemma~\ref{lem:coercivity} and the upper bound from Lemma~\ref{lem:controllability}. The definition of $K:=\bar{C}{\epsilon}/({\epsilon-1})$ with $\bar{C}$ from Lemma \ref{lem:controllability} then leads to \eqref{eq:prop_TP_1}, which finishes this proof.
\end{proof}

\end{document}
